\subsection{Differentiator}\label{sec:amp_differentiator}

\subsubsection{Overview}
An inverting amplifier with a capacitor \(C\) at the input and \(G_2\) in
feedback. The output is the negative time derivative of the input.

\subsubsection{Schematic}
\begin{circuitfig}
  \draw (0,0) node[op amp](A){};
  \draw (A.-) -- ++(-0.6,0) coordinate(n) node[circ]{} to[capacitor, l_={$C$}] ++(-1.6,0) to[european resistor, l_={$G_s$ (optional)}] ++(-2,0) node[circ]{} node[left]{$V_{in}$};
  \draw (n) |- (-0.5,1.6) to[european resistor, l={$G_2$}] (1.6,1.6) -| (A.out);
  \draw (A.+) -- ++(-0.6,0) node[ground]{};
  \draw (A.out) to[short, -o] ++(0.8,0) node[right]{$V_{out}$};
\end{circuitfig}

\subsubsection{Transfer Function}
\textbf{Small Signal, Frequency Domain, AC Input Analysis (Q4)}, \(A = A(s)\),
\(G_s \to \infty\): \(V_- = -\frac{V_{out}}{A}\). Applying KCL at \(V_-\):
\begin{align*}
  (V_- - V_{in}) sC + (V_- - V_{out}) G_2 &= 0 \\
  V_- (sC + G_2) - V_{out} G_2 &= V_{in} sC \\
  -\frac{V_{out}}{A}(sC + G_2) - V_{out} G_2 &= V_{in} sC \hspace{1cm} \text{(Substituting \(V_- = -\frac{V_{out}}{A}\))} \\
  -V_{out}\left(\frac{sC + G_2}{A} + G_2\right) &= V_{in} sC \\
  H(s) &= -\frac{sC}{G_2 + \frac{sC + G_2}{A}} \xrightarrow{A \to \infty} -\frac{sC}{G_2} = -s R_2 C
\end{align*}
\textbf{Small Signal, Time Domain, AC Input Analysis (Q2)}: multiplication by
\(s\) is differentiation (zero initial conditions):
\[
  V_{out}(t) = -R_2 C \frac{dV_{in}}{dt}
\]

\subsubsection{Analysis}
\begin{enumerate}
  \item \textbf{Stability}: with \(V_{in} = 0\), KCL at \(V_-\):
    \begin{align*}
      V_- sC + (V_- - V_{out}) G_2 &= 0 \\
      V_- &= \frac{G_2}{G_2 + sC} V_{out} \hspace{1cm} \Rightarrow \hspace{1cm}
      F(s) = \frac{G_2}{G_2 + sC} = \frac{1}{1 + s R_2 C} \hspace{0.5cm} [\div G_2]
    \end{align*}
    The feedback factor carries a pole. With \(A(s) = \frac{A_0 p_1}{s + p_1}\)
    the closed loop poles are the roots of \(1 + A F = 0\):
    \begin{align*}
      1 + \frac{A_0 p_1}{(s + p_1)(1 + s R_2 C)} &= 0 \\
      (s + p_1)(1 + s R_2 C) + A_0 p_1 &= 0 \hspace{1cm} [\times (s + p_1)(1 + s R_2 C)] \\
      s + s^2 R_2 C + p_1 + p_1 R_2 C s + A_0 p_1 &= 0 \\
      R_2 C s^2 + (1 + p_1 R_2 C) s + p_1 (1 + A_0) &= 0 \\
      s^2 + \frac{1 + p_1 R_2 C}{R_2 C} s + \frac{p_1 (1 + A_0)}{R_2 C} &= 0 \hspace{1cm} [\div R_2 C]
    \end{align*}
    Comparing with \(s^2 + \frac{\omega_n}{Q} s + \omega_n^2\)
    (Section~\ref{sec:theory_cs_second_order}), with \(p_1 (1 + A_0) \approx p_1 A_0 = 2\pi GB\):
    \begin{align*}
      \omega_n^2 &= \frac{p_1 (1 + A_0)}{R_2 C} \approx \frac{2\pi GB}{R_2 C} \hspace{1cm} \Rightarrow \hspace{1cm} \omega_n = \sqrt{\frac{2\pi GB}{R_2 C}} \\
      \frac{\omega_n}{Q} &= \frac{1 + p_1 R_2 C}{R_2 C} \hspace{1cm} \Rightarrow \hspace{1cm}
      Q = \frac{\omega_n R_2 C}{1 + p_1 R_2 C} = \frac{\sqrt{\frac{2\pi GB}{R_2 C}} R_2 C}{1 + p_1 R_2 C} = \frac{\sqrt{2\pi GB R_2 C}}{1 + p_1 R_2 C}
    \end{align*}
    E.g. \(GB = 1\,\)MHz, \(R_2 C = 1\,\)ms, \(p_1 = 2\pi \cdot 10\,\)rad/s:
    \begin{align*}
      2\pi GB R_2 C &= 2\pi \cdot 10^6 \cdot 10^{-3} = 6283 \hspace{1cm} \sqrt{6283} = 79.27 \\
      p_1 R_2 C &= 2\pi \cdot 10 \cdot 10^{-3} = 0.0628 \hspace{1cm}
      Q = \frac{79.27}{1 + 0.0628} = \frac{79.27}{1.0628} = 74.6 \\
      \omega_n &= \sqrt{\frac{2\pi \cdot 10^6}{10^{-3}}} = \sqrt{6.283 \times 10^{9}} = 7.927 \times 10^{4}\,rad/s \hspace{1cm}
      f_n = \frac{\omega_n}{2\pi} = 12.6\,kHz
    \end{align*}
    The circuit rings at \(12.6\,\)kHz. It also amplifies high frequency
    noise (\(|H(j\omega)| = \omega R_2 C \propto f\)).
  \item \textbf{Practical Form}: a series \(G_s\) with \(C\) has admittance
    \(\frac{1}{\frac{1}{sC} + \frac{1}{G_s}} = \frac{sC G_s}{sC + G_s}\); with the ideal op-amp
    (\(V_- = 0\)):
    \begin{align*}
      (0 - V_{in}) \frac{sC G_s}{sC + G_s} + (0 - V_{out}) G_2 &= 0 \\
      H(s) &= -\frac{\frac{sC G_s}{sC + G_s}}{G_2} = -\frac{sC G_s R_2}{sC + G_s} = -\frac{s R_2 C}{1 + s R_s C} \hspace{1cm} [\div G_s]
    \end{align*}
    For \(f \ll \frac{1}{2\pi R_s C}\), \(H \approx -s R_2 C\) (differentiates); for
    \(f \gg \frac{1}{2\pi R_s C}\), \(H \approx -\frac{s R_2 C}{s R_s C} = -\frac{R_2}{R_s}\) (flat).
\end{enumerate}
